% ================================================================
% ################################################################
% ================================================================
\section{The Decision Points in Detail}
\label{sec:formal}

This appendix expands on the decision points of \cref{tab:divergence}: the
choices an evaluator must fix but the STL formula does not determine. We
follow the three groups of the table, discussing each decision point in
turn. The signal interpretation, which bundles the several choices of
\cref{tab:siginterp}, is treated choice by choice, and for the two headline
choices of \cref{sec:discrepancies}---the boundary rule and the
terminal-step rule---we show how each produces the divergence reported
there. We first recall the standard robustness semantics that every
evaluator approximates.

% ------------------------------------------------
\subsection{Standard Robustness}

Fix a signal $\sigma$ in the sense of \cref{def:signal} and a formula
$\varphi \in \mathbf{Fml}$ as in \cref{def:stl}. The standard quantitative
(robustness) semantics
$\rho(\varphi, \sigma, t) \in \mathbb{R} \cup \{\pm\infty\}$ is defined
inductively, following \cite{fainekos2009robustness,donze2010robust}:
\begin{align*}
  \rho(\top,\sigma,t) &= +\infty \\
  \rho(f(\vec{x}) \geq 0,\sigma,t) &= f(\sigma_{\vec{x}}(t)) \\
  \rho(\neg\varphi,\sigma,t) &= -\rho(\varphi,\sigma,t) \\
  \rho(\varphi_1 \vee \varphi_2, \sigma, t) &= \max(\rho(\varphi_1,\sigma,t),\,\rho(\varphi_2,\sigma,t)) \\
  \rho(\varphi_1 \wedge \varphi_2, \sigma, t) &= \min(\rho(\varphi_1,\sigma,t),\,\rho(\varphi_2,\sigma,t)) \\
  \rho(\varphi_1 \mathbin{U_I} \varphi_2, \sigma, t) &= \sup_{t' \in t + I} \min\Bigl(\rho(\varphi_2,\sigma,t'),\, \inf_{t'' \in [t,\,t')} \rho(\varphi_1,\sigma,t'')\Bigr) \\
  \rho(\Diamond_I\varphi, \sigma, t) &= \sup_{t' \in t + I} \rho(\varphi,\sigma,t') \\
  \rho(\Box_I\varphi, \sigma, t) &= \inf_{t' \in t + I} \rho(\varphi,\sigma,t'),
\end{align*}
where $t + I = \{t + s \mid s \in I\}$ and $f(\sigma_{\vec{x}}(t))$ applies
$f$ to the values of the variables $\vec{x}$ at time $t$. The release
operator is defined by duality,
$\varphi_1 \mathbin{R_I} \varphi_2 :\equiv
 \neg(\neg\varphi_1 \mathbin{U_I} \neg\varphi_2)$,
and the constant $\bot :\equiv \neg\top$ has
$\rho(\bot,\sigma,t) = -\infty$.

Read literally, this definition assumes an idealized signal: continuous,
infinitely long, and known at every instant. An evaluator never has that. It
has a finite list of samples $\sigma(t_0), \dots, \sigma(t_n)$, and to compute the
$\sup$, $\inf$, and $\min$ above it must decide what the signal does at the
times that fall between, before, and after those samples. Those decisions
form the signal interpretation discussed in \cref{sec:app-open}; the
remaining decision points arise from corner cases of the definition itself
(\cref{sec:app-corner}) and from extensions that applications demand beyond
it (\cref{sec:app-ext}).

% ------------------------------------------------
\subsection{Fragments and Corner Cases of the Standard Definition}
\label{sec:app-corner}

\paragraph{Support for until and release.}
Few tools implement the full syntax of \cref{def:stl}. The until operator,
with its nested $\inf$, is harder to implement and to vectorize than $\Box$
and $\Diamond$, and release is often dropped together with it: some tools
support only the box/diamond fragment, others support until under
restricted interval forms. The effect is on coverage: a specification
accepted by one tool is rejected by another before any robustness is
computed.

\paragraph{Open vs.\ closed interval endpoints.}
\cref{def:stl} admits open and half-open intervals such as $(1,3]$, but many
front ends parse only closed ones. Where open intervals are accepted, their
reading depends on the signal interpretation: under a continuous interpolant
the $\inf$ over an open window coincides with that over its closure, so the
difference is invisible, while under a discrete reading excluding an
endpoint drops a sample and can change the aggregated value.

\paragraph{Robustness of $\top$ and $\bot$.}
The standard semantics assigns $\rho(\top,\sigma,t) = +\infty$ and
$\rho(\bot,\sigma,t) = -\infty$. Tools differ in whether they admit infinite
values at all: some propagate genuine infinities through the $\min$/$\max$
aggregation, others substitute large finite sentinels, and a sentinel leaks
into reported values once it is aggregated, scaled, or differentiated.

% ------------------------------------------------
\subsection{Choices the Definition Leaves Open}
\label{sec:app-open}

The central decision point in this group is the \emph{signal
interpretation}: the bundle of choices, summarized in \cref{tab:siginterp},
by which the finite timed state sequence an evaluator receives is read as
the signal $\sigma$ that the semantics quantifies over. We discuss its
choices in turn, then the remaining decision points of the group.

\paragraph{Signal model.}
The signal model decides how a sampled trace is read as a function of
continuous time. Under a \emph{piecewise-linear} dense-time model (Breach, and
\tool's native backend) the value between two samples is their linear
interpolant, so a predicate such as $x \geq 0$ can change truth between sample
points. A \emph{piecewise-constant} model holds each sample until the next. A
\emph{discrete-time} model (RTAMT's default) drops the notion of between-sample
time altogether: the trace is an indexed sequence and temporal operators count
indices, not durations. The same formula and the same samples can therefore
denote different signals before evaluation even begins.

\paragraph{Window alignment.}
Even once the signal model is fixed, a temporal window $[t+a,\,t+b]$ has
endpoints that need not coincide with samples. The window-alignment rule
decides the value taken there. A dense-time backend evaluates the interpolant
at the exact endpoint; a discrete backend snaps to the nearest index or
ignores the fractional part. This choice is invisible on uniformly sampled
traces whose windows align with samples, which is why it is easy to overlook,
and why divergence appears only on alignment-sensitive inputs.

\paragraph{Boundary rule.}
When a window $[t+a,\,t+b]$ reaches past the end of the trace at $T$, the
$\sup$ or $\inf$ ranges in part over times where no sample exists. The boundary
rule supplies those values. Under \emph{clamp} (\tool, Breach) the last
sample is treated as persisting, so an out-of-range query returns $\sigma(T)$. Under
a \emph{pessimistic} rule (RTAMT discrete) the out-of-range part contributes
nothing: the window is effectively empty, so a $\Diamond$ becomes $-\infty$
(the $\sup$ of an empty set) and a $\Box$ becomes $+\infty$.
This is the boundary-rule row of the divergence block in
\cref{tab:divergence-matrix}. For
$\varphi = \Diamond_{[2,4]}(x \geq 0)$ evaluated at $t = 1$ on a trace that
ends at $T = 2$ with $x(T) = 3$, the window $[3,5]$ lies entirely beyond the
trace. Clamp sees the persisting value $3$ and reports $+3$ (satisfied);
the pessimistic rule sees an empty window and reports $-\infty$ (violated).
Same formula, same samples, same time point; opposite verdicts.

\paragraph{Terminal-step rule.}
At the final sample $t_N$ a tool may report the robustness it computed there,
or it may overwrite it with the value from the penultimate sample $t_{N-1}$.
Breach does the latter by default; RTAMT does the former. The two
agree whenever the final and penultimate robustness coincide, so the choice is
silent on most traces. It is not silent when they differ.
This is the terminal-step row of the divergence block in
\cref{tab:divergence-matrix}. For
$\varphi = (x \geq 0) \wedge (y \geq 0)$ on trace~A with $x = y = [5, 5, -3]$,
the penultimate robustness is $+5$ and the final is $-3$. The standard rule
reports $-3$; the extension rule copies the penultimate $+5$. Against a
trace~B that holds $+0.1$ throughout, this flips which trace looks better,
reversing a ranking that a falsification search would act on.

\paragraph{Range of expressible atomic predicates.}
\cref{def:stl} permits an arbitrary function symbol $f$, but front ends
restrict what is expressible: to threshold comparisons on single variables,
to linear arithmetic, or to a fixed expression grammar. The effect is on
coverage.

\paragraph{Robustness of atomic predicates.}
Even for an accepted predicate, the evaluator must decide how it is scored:
the robustness of $f(\vec{x}) \geq 0$ may be the raw signed margin
$f(\vec{x})$ or a Euclidean distance to the satisfaction
set~\cite{fainekos2009robustness}, and the two disagree once a predicate
involves several variables.

% ------------------------------------------------
\subsection{Domain-Specific Extensions}
\label{sec:app-ext}

The third group extends the definition rather than reads it: applications
genuinely demand more than \cref{def:stl} provides, and each extension is a
decision point because adopting a tool silently adopts its extensions along
with their semantics.

\paragraph{Equality predicates.}
The equality $x = 0$ is absent from \cref{def:stl} yet common in practical
specifications. It is satisfied only at the single value $x = 0$ and admits
no single standard robustness; a common choice is $-|x|$, which is never
positive, so even a satisfied equality reports non-positive robustness.
Breach instead scores equality by a fixed-magnitude sign convention
($\pm 10^4$), so the reported value carries no metric information at
all---the equality row of \cref{tab:divergence-matrix} shows the resulting
cross-tool divergence on inputs every tool accepts.

\paragraph{Past-time operators.}
Operators such as once and historically mirror the temporal modalities into
the past. They change the direction in which an evaluator walks the trace,
and tools either provide them (RTAMT) or not. The effect is on coverage.

\paragraph{Differentiable robustness.}
Gradient-based learning replaces the $\min$/$\max$ aggregations by smooth
approximations~\cite{levenberg2020stlcg,pant2017smooth}. The smooth value
deviates from the standard robustness by construction; the decision point is
whether the deviation is an exposed, tunable parameter or an implicit
consequence of the tool choice.

\paragraph{Online evaluation.}
Online monitoring evaluates a growing prefix, so verdicts must be issued
before the trace is complete~\cite{deshmukh2017online}. This forces a
treatment of missing future information---three-valued verdicts, robustness
intervals---that offline evaluation never confronts.

% ------------------------------------------------
\subsection{Why ``Implicit''}

None of these choices is written in the STL formula; each is fixed by the
evaluator. Choosing a tool therefore silently chooses a combination of them,
and two tools whose combinations differ can report different robustness, or
even opposite verdicts, on identical input, as the boundary and terminal-step
examples show. Nor is \cref{tab:divergence} a closed taxonomy: the point is
not completeness but that such choices exist, are fixed by the evaluator, and
should be made explicit. The contribution of \tool is not to declare one
combination correct, but to turn the combination into an explicit, switchable
argument rather than a property baked into the tool.

% ================================================================
% ################################################################
% ================================================================
\section{The Trace Interface}
\label{sec:app-trace}

This appendix records how the inspectability used in
\cref{sec:discrepancies} is realized; the rationale is in
\cref{sec:design}. The facts below follow the implementation.

Every backend's \texttt{evaluate(formula, signal)} returns a result
object satisfying one protocol: \texttt{robustness} (the $(N,T)$ output
array; \texttt{robustness(...)}\ in the public API is shorthand for it),
a \texttt{has\_trace} flag, and two trace methods,
\texttt{trace\_for(node)} and \texttt{traced\_nodes()}. During lowering,
the backend records for every DAG operation it emits the sub-formula
node it originates from; the operation that produces a sub-formula's
value is that sub-formula's \emph{representative}. The executor
evaluates the DAG in topological order and retains every operation's
output, so the result carries all intermediate robustness signals, not
only the root. \texttt{trace\_for(sub)} returns the representative
output for the sub-formula \texttt{sub}; \texttt{traced\_nodes()}
enumerates the sub-formulas with observable traces.

\begin{pycode}
phi = parse("G[0,5]((x >= 0) and F[0,2](y >= 0))")
res = evaluate(phi, sig, backend="native")
sub = phi.children[0].children[1]   # the F[0,2] sub-formula
rho_sub = res.trace_for(sub)        # its (N, T) robustness trace
\end{pycode}

\noindent
Three details matter for the experiments of \cref{sec:discrepancies}.
First, \texttt{trace\_for} is keyed by the identity of the sub-formula
node, so one parsed formula evaluated under two backends yields
directly comparable per-sub-formula traces; the localizer of
\cref{sec:discrepancies} is a post-order walk over
\texttt{node.children} that diffs \texttt{trace\_for} at each node.
Second, backend post-processing participates in the trace: the
Breach-compatible backend's terminal-step extension replaces the
top-level output in the result, which is why the localizer reports it
(the ``\texttt{PointwiseMin + terminal-step extension}'' line). Third,
trace support is per backend: the Rust SIMD executor declares
\texttt{has\_trace = False} and raises on the trace methods---the
per-operation retention is exactly the cost its scaling column drops.

% ================================================================
% ################################################################
% ================================================================
\section{The Tool Landscape in Detail}
\label{sec:app-landscape}

\cref{tab:landscape-detailed} compares representative tools along the axes
on which their evaluators actually differ, expanding the summary in
\cref{sec:background}.

\begin{table}[ht]
\centering
\caption{How representative tools differ along the axes that shape their
evaluators.}
\label{tab:landscape-detailed}
\resizebox{\textwidth}{!}{%
\begin{tabular}{lllllll}
\hline
Tool & Interface & Signal model & Online mode & Primary target & Auto-grad & Stack \\
\hline
Breach~\cite{donze2010breach,donze2013efficient} & API + DSL & piecewise-linear   & offline & falsification           & no  & MATLAB, mature \\
RTAMT~\cite{nickovic2020rtamt,nickovic2024rtamt} & DSL + API & discrete (default) / dense & delayed & runtime verification    & no  & Python, heavy \\
MoonLight~\cite{bartocci2023moonlight}           & DSL + API & spatial + temporal & partial & spatial RV              & no  & Java / CLI / Python \\
STLCG++~\cite{kapoor2025stlcgpp}                 & API       & discrete (tensor)  & none    & differentiable learning & yes & PyTorch \\
% \tool     & DSL + API & piecewise-linear   & offline & evaluation              & no  & Python, light \\
\hline
\end{tabular}%
}
\end{table}

% ================================================================
\section{Experimental Details}
\label{sec:app-eval}
% ================================================================

This appendix records the protocols and full data behind the findings
reported in \cref{sec:evaluation}.

\paragraph{Case set and comparison protocol.}
The baseline block comprises 16 operator-coverage cases drawn from the
regression suite; the divergence block adds five cases, each targeting one
decision point of \cref{tab:divergence,tab:siginterp}. Each tool runs in
its native environment (versions as in \cref{tab:divergence-matrix}).
Outputs are normalized to per-timestep robustness vectors and compared
pairwise at absolute tolerance $10^{-6}$ on mutually defined timesteps.
STLCG++ runs with exact $\min$/$\max$ aggregation---its default
configuration; smooth aggregation is opt-in---so its column reflects its
discrete-time semantics rather than smoothing, which remains a
coverage-level decision point of \cref{tab:divergence} not exercised here.

\paragraph{Compatibility validation.}
Every tool column of \cref{tab:divergence-matrix} doubles as a validation
check: on each case, the named compatibility backend must reproduce its
tool's output per timestep at $10^{-6}$---21 cases $\times$ 3 tools, hence
63 case--tool cells. The check passes on all 59 cells that carry values;
three further cells match on coverage (backend and tool agree the case is
inexpressible). The one mismatch is deliberate: where real RTAMT silently
reinterprets a non-uniform grid as uniformly indexed samples, the
RTAMT-compatible backend rejects the input instead of reproducing the
reinterpretation. The backends implement each tool's \emph{documented}
rules; the reproduction check verifies the implementations rather than
defining them---and it is nontrivial: on the first live Breach run it
exposed two fidelity bugs in our own Breach-compatible backend (the
equality convention and the window-endpoint reduction), both fixed and
re-verified before the numbers reported here were produced.

\paragraph{Localization procedure.}
On the validated columns, the localizer exploits that \tool backends share
one formula representation with per-node traces: it walks the shared
syntax tree in post-order and reports the \emph{minimal divergent
node}---the lowest node whose traces disagree while all its children
agree---together with each backend's primitive operation at that node. On
the two headline rows of \cref{tab:divergence-matrix}:
\begin{lstlisting}[language={},keywordstyle={},basicstyle=\footnotesize\ttfamily]
first divergent op at `F[2,4]`: WindowMax[2,4] vs
    DiscreteWindowMax[2,4] (t-index 1) -> boundary rule
first divergent op at `and`: PointwiseMin vs PointwiseMin
    + terminal-step extension (t-index 2) -> terminal-step rule
\end{lstlisting}
Backend pairs whose lowerings differ structurally are compared at the
shared syntax-tree level rather than between lowered representations;
predicate-internal arithmetic is treated as a leaf, since no decision
point lives inside it.

\paragraph{Trace-family construction.}
A single constructed trace could sit on a knife edge, so each headline
case is widened to a swept family of 201 traces whose final sample varies
over $[-5,5]$, spanning regions of agreement and disagreement; no rule is
taken as the reference, and the rates are symmetric between the rules by
construction. For the boundary-rule family ($\Diamond_{[2,4]}(x \ge 0)$,
window past the trace end), the clamping and pessimistic rules return
opposite verdicts at the divergence timestep on exactly those traces whose
final sample is non-negative (101 of 201). The terminal-step family is a
conjunction violated only at the final sample; the as-is and
extend-penultimate rules flip the final verdict on 100 of 201 traces and
rank the trace oppositely against a fixed reference trace with an interior
violation---the ranking signal a falsification search acts on---on 80 of
201. These rates measure the width of the disagreement region within
each controlled family---the divergences are not measure-zero
coincidences---not a frequency claim about any particular workload.

\paragraph{Falsification setup.}
The system is a closed-form speed model $\dot v = K u(t) - c\,v$ with
piecewise-constant throttle on four segments (a bounded $[0,1]^4$ search
space; traces sampled exactly, no integrator); the specification is
$(v \le 8.8) \wedge \Diamond_{[0.5,1.5]}(v \le 25)$, whose second conjunct
never fails but carries the monitor horizon past the end of the trace. A
fixed $(1{+}1)$-ES with random restarts minimizes the monitor-style
objective $\min_t \rho(t)$; only the objective semantics varies across
arms, and every claimed falsification is validated by an exact closed-form
check of the system itself, so no backend acts as the verdict oracle. The
configuration, seeds, and budget (50 seeded runs, 300 simulations each)
were frozen before the runs.

Per-arm results behind \cref{fig:falsification}: under the clamping
boundary rule the search falsifies in 50/50 runs (median 26.5
simulations); adding the terminal-step extension masks violations that
appear only at the final sample, costing one run and roughly $1.5\times$
the simulations (49/50, median 40). Under the pessimistic rule the
out-of-range window makes the objective $-\infty$ for every input, so the
search degenerates to blind sampling: 10/50 runs falsify (median 214), and
14{,}027 candidate falsifications claimed across the campaign are rejected
by the exact system-level check. The search runs to completion in every
arm; the disagreement between the arms is silent.

\paragraph{Variant accounting.}
LOC in \cref{tab:variation} counts variant-specific code only: blank,
comment-only, and docstring lines are excluded, shared infrastructure is
excluded, and Rust is counted separately. The native backend is the
baseline and has no row; the Rust SIMD executor ships as an optional
native extension alongside the five backends. The full regression suite
comprises 181 tests and passes with all variants installed.

\paragraph{Full timing data.}
\cref{tab:scaling} tabulates the measurements plotted in
\cref{fig:scaling}.

\begin{table}[ht]
\centering
\caption{Mean evaluation time (ms) $\pm$ std over a batch of $N$ traces of
length $T$; identical workload
$\Box_{[0,5]}((x>0)\wedge\Diamond_{[0,2]}(y>0))$ for every column. This is
the data plotted in \cref{fig:scaling}. $\dagger$: RTAMT evaluates one
trace per call, so batch times are $N \times$ the measured per-trace time
(linearity checked at $N \in \{1,8\}$). OOM: window materialization
exceeds memory on the 32\,GB host.}
\label{tab:scaling}
% AUTO-GENERATED from extra/artifacts/rv2026-eval/e4/scaling_merged.json
% by extra/experiments/e4_merge_scaling.py (tidystl repo); do not hand-edit.
% Times in ms, mean +- std over repeats, identical workload G[0,5]((x>0) and F[0,2](y>0)).
% dagger: N x measured per-trace time (RTAMT is single-trace; its
% batching model is a loop, linearity measured at N in {1,8}).
% OOM: window materialization exceeds memory on the 31 GiB host.
\begin{tabular}{rr r r r r r r}
\hline
$N$ & $T$ & Native & Breach-c. & RTAMT-c. & SIMD & RTAMT & STLCG++ \\
\hline
   1 &   101 & $1.30 \pm 0.08$ & $0.98 \pm 0.06$ & $0.57 \pm 0.06$ & $0.02 \pm 0.01$ & $0.40 \pm 0.01$ & $0.31 \pm 0.05$ \\
   1 &  1001 & $12.15 \pm 0.19$ & $9.34 \pm 0.28$ & $4.58 \pm 0.08$ & $0.19 \pm 0.01$ & $3.74 \pm 0.28$ & $33.71 \pm 1.31$ \\
  64 &   501 & $11.75 \pm 0.25$ & $7.94 \pm 0.15$ & $5.24 \pm 0.20$ & $2.44 \pm 0.41$ & $118^\dagger$ & $224 \pm 3$ \\
  64 &  5001 & $358 \pm 11$ & $167 \pm 6$ & $136 \pm 2$ & $31.20 \pm 1.89$ & $1195^\dagger$ & OOM \\
 256 &   501 & $26.19 \pm 1.54$ & $21.42 \pm 0.65$ & $15.97 \pm 0.37$ & $3.09 \pm 0.50$ & $471^\dagger$ & $934 \pm 44$ \\
 256 &  5001 & $1377 \pm 155$ & $569 \pm 9$ & $515 \pm 23$ & $115 \pm 2$ & $4781^\dagger$ & OOM \\
\hline
\end{tabular}%

\end{table}
